%% LyX 2.5.1 created this file.  For more info, see https://www.lyx.org/.
%% Do not edit unless you really know what you are doing.
\documentclass[english,hebrew]{article}
\usepackage[HE8,T1]{fontenc}
\usepackage[utf8]{inputenc}
\usepackage{amssymb}
\usepackage{babel}
\begin{document}
\title{מבוא לטופולוגיה קבוצתית: קבוצות סגורות במרחבים מטריים}
\maketitle
\begin{description}
\item [{קטגוריות:}] טופולוגיה
\item [{תגים:}] מרחבים מטריים, מרחבים טופולוגיים
\item [{מזהה:}] \L{closed\_sets\_in\_metric\_spaces}
\end{description}
בפוסט הקודם דיברתי על המושג של \textbf{קבוצה פתוחה}. במרחב מטרי, זו
הייתה מין הכללה של כדור פתוח: קבוצה \L{$A$} היא פתוחה אם לכל \L{$a\in A$}
קיים כדור פתוח \L{$B\left(a,r\right)$} עם רדיוס \L{$r>0$} כך ש-\L{$B\left(a,r\right)\subseteq A$}.
בעזרת המושג הזה הגדרתי \textbf{סביבה} של נקודה בתור קבוצה פתוחה שמכילה
אותה, ובעזרת מושג הסביבה הגדרתי מחדש מה זה גבול של סדרה: \L{$\lim_{n\to\infty}a_{n}=a$}
אם לכל סביבה \L{$U$} של \L{$a$}, קיים \L{$N$} כך שלכל \L{$n>N$}
מתקיים ש-\L{$a_{n}\in U$}.

עכשיו בואו נניח ש-\L{$A$} היא קבוצה כלשהי, לאו דווקא פתוחה. נסתכל
על כל הסדרות של איברים מתוך \L{$A$}. מן הסתם חלק נכבד מהסדרות הללו
לא יתכנסו, אבל מה עם אלו שכן? האם אפשר להגיד משהו על הקשר בין הגבולות
שלהן ל-\L{$A$}? ראשית, כל איבר \L{$a\in A$} הוא בעצמו גבול של סדרה
של איברים מתוך \L{$a$} - הסדרה הקבועה \L{$a,a,a,\ldots$}. אבל בואו
נניח עכשיו שנקודה \L{$a\notin A$} היא גם כן גבול של איברים מ-\L{$A$}.
זה אומר, על פי הגדרת הגבול שנתתי, ש\textbf{כל} סביבה של \L{$a$} מכילה
אינסוף איברים מתוך הסדרה - כל האיברים החל ממקום כלשהו. האם זה אומר
שכל סביבה של \L{$a$} מכילה אינסוף איברים מתוך \L{$A$}? ובכן, לא,
כי מבין אינסוף האיברים בסדרה ייתכן שיש רק מספר סופי של איברים \textbf{שונים}
והיתר חוזרים על עצמם )לפני שניה ראינו סדרה אינסופית שמורכבת כולה מאיבר
אחד(. אז מה שאני כן יכול להגיד הוא שבכל סביבה של \L{$a$} יש לפחות
איבר אחד של \L{$A$}. לתכונה הזו יש שם: \textbf{נקודת הצטברות}. פורמלית,
\L{$x$} היא נקודת הצטברות של \L{$A$} אם בכל סביבה של \L{$x$} יש
איבר \L{$a\ne x$} כך ש-\L{$a\in A$} )כלומר, אם \L{$x\in A$} זה
עדיין לא מבטיח אוטומטית שהוא נקודת הצטברות(.

אוקיי, אז ראינו שאם \L{$a\notin A$} הוא \textbf{גבול של סדרה} אז
\L{$a$} הוא \textbf{נקודת הצטברות} של \L{$A$}. האם ההפך נכון? האם
כל נקודת הצטברות של \L{$A$} היא גבול של סדרה של איברים מ-\L{$A$}?
זה תרגיל ששווה לנסות לחשוב עליו רגע עצמאית כדי לראות מה עובד ומה לא
עובד. לעת עתה, בואו נזכור שאנחנו עדיין מדברים בהקשר של \textbf{מרחבים
מטריים} כי לא הגדרתי מרחבים מסוג אחר, וננסה להבין איך אפשר להוכיח
את זה עבורם. זה לא כזה קשה, פשוט משתמשים בתעלול שגור באינפי. בהינתן
\L{$a$} נגדיר סדרה של \textbf{כדורים פתוחים} סביב \L{$a$}, הסדרה
\L{$A_{n}=B\left(a,\frac{1}{n}\right)$}. הכדורים הללו הם מרדיוס שהולך
וקטן, ולכן יש לנו סדרת הכלות \L{$A_{1}\supseteq A_{2}\supseteq A_{3}\supseteq\ldots$}.

עכשיו, מכיוון ש-\L{$a$} היא נקודת הצטברות של \L{$A$}, בכל סביבה
שלה יש איבר של \L{$A$}. כדור פתוח הוא סביבה של \L{$a$}, ולכן לכל
\L{$A_{n}$} כזה קיים \L{$a_{n}\in A$} כך ש-\L{$a_{n}\in A_{n}$}.
קיבלנו את הסדרה שלנו - עדיין צריך להראות שהיא מתכנסת אל \L{$a$}.

אם כן, בואו ניקח סביבה \textbf{כלשהי} של \L{$a$}, נסמן אותה \L{$U$}.
סביבה היא קבוצה פתוחה, ולכן מההגדרה של קבוצה פתוחה שנתתי במשפט הפותח
בפוסט הזה, קיים כדור פתוח \L{$B\left(a,r\right)$} עם \L{$r>0$} כך
ש-\L{$B\left(a,r\right)\subseteq A$}. עכשיו ניקח \L{$N$} כך ש-\L{$\frac{1}{N}<r$},
אז יתקיים \L{$B\left(a,\frac{1}{N}\right)\subseteq B\left(a,r\right)\subseteq U$}.
עכשיו, לכל איבר \L{$a_{n}$} בסדרה שמקיים \L{$n\ge N$} יתקיים

\L{$a_{n}\in B\left(a,\frac{1}{n}\right)\subseteq B\left(a,\frac{1}{N}\right)\subseteq U$}

וזה מה שרצינו להראות.

השתמשנו פה בכל מני דברים שקשורים למרחבים מטריים, בפרט תכונות הכלה
של כדורים; השתמשנו כאן גם משהו שנקרא \textbf{התכונה הארכימדית} של
המספרים, שהיא מה שנותן לנו את זה שקיים \L{$N$} עבורו \L{$\frac{1}{N}<r$}.
בקיצור, לא מעט מתמטיקה קונקרטית ו\textquotedblright לא אבסטרקטית\textquotedblright .
מה התכונה האבסטרקטית שמתחבאת כאן והייתה נותנת לנו את ההוכחה, והאם
זה משהו שמתקיים בכל מרחב טופולוגי?

ובכן, לא, התכונה הזו \textbf{לא} מתקיימת בכל מרחב טופולוגי. כשנגיע
להגדרה של מרחבים טופולוגיים נראה דוגמאות למרחבים שבהן יש נקודות הצטברות
שלא יכולות להיות הגבולות של סדרות - אינטואיטיבית בגלל ש\textquotedblright צריך
מספר לא בן מניה של נקודות\textquotedblright{} כדי להשתכנע שהן נקודות
הצטברות. נשמע מוזר? יופי! עכשיו ברור אילו עולמות מוזרים נחשפים ככל
שאנחנו מבצעים אבסטרקציה להנחות שלנו. אבל מה התכונה ש\textbf{כן} מבטיחה
שכל נקודת הצטברות תהיה גבול?

בהוכחה עם מרחבים מטריים שראינו, השתמשתי בתכונות של מרחבים מטריים כדי
לבנות סדרה \L{$A_{1},A_{2},A_{3},\ldots$} של סביבות של \L{$a$}.
הסדרה הזו הייתה בעלת שתי תכונות מועילות:
\begin{enumerate}
\item לכל סביבה \L{$U$} של \L{$a$} היה קיים \L{$N$} כך ש-\L{$A_{N}\subseteq U$}.
\item התקיים \L{$A_{1}\supseteq A_{2}\supseteq A_{3}\supseteq\ldots$}.
\end{enumerate}
השילוב של שתי התכונות הללו הוא מה שנתן לנו את ההוכחה. אם כן, האם צריך
לדרוש את שתיהן? ובכן, לא; אפשר להסתפק בטיפה פחות - התכונה הראשונה\textbf{
}תספיק.

בואו נראה איך זה עובד. אני מניח שיש לי סדרה \L{$A_{1},A_{2},A_{3},\ldots$}
של סביבות של \L{$a$} בעלת התכונה שלכל סביבה \L{$U$} של \L{$a$},
קיים \L{$N$} כך ש-\L{$A_{N}\subseteq U$}. הסביבות הללו \textbf{לא}
מקיימות בהכרח את התכונה \L{$A_{1}\supseteq A_{2}\supseteq A_{3}\supseteq\ldots$}.
אז אני הולך לשנות אותן קצת. אני אגדיר סדרה \textbf{חדשה} של סביבות,
\L{$C_{1},C_{2},C_{3},\ldots$} שתקיים את שתי התכונות. לשם כך אני
אגדיר \L{$C_{1}=A_{1}$} ומעכשיו באינדוקציה, \L{$C_{n}=C_{n-1}\cap A_{n}$}.
במילים אחרות, \L{$C_{n}=A_{1}\cap A_{2}\ldots\cap A_{n}$}. באופן
מובהק זה הולך לתת לנו את תכונת ההכלה \L{$C_{1}\supseteq C_{2}\supseteq C_{3}\supseteq\ldots$}
פשוט כי כל קבוצה מתקבלת מהקודמת על ידי חיתוך של הקודמת עם עוד איזה
משהו.

האם גם תכונה {\beginL 1\endL} מתקיימת? בוודאי. כי אם \L{$U$} היא
סביבה של \L{$a$} וקיים \L{$N$} כך ש-\L{$A_{N}\subseteq U$} אז מכך
ש-\L{$C_{N}=C_{N-1}\cap A_{N}\subseteq A_{n}$} נקבל ש-\L{$C_{N}\subseteq A_{N}$}.

מה עוד אולי חסר לנו? חסר לנו ש-\L{$C_{n}$}-ים הללו הם \textbf{סביבות}
של \L{$a$}. כלומר, ברור ש-\L{$a\in C_{n}$} לכל \L{$n$} כי קל להוכיח
את זה באינדוקציה תוך הסתמכות על כך ש-\L{$a\in A_{n}$} לכל \L{$n$}.
אבל למה הקבוצות \L{$C_{n}$} הללו הן \textbf{פתוחות}? כדי שמשהו ייקרא
\textquotedblleft סביבה\textquotedblright{} הוא צריך להיות קבוצה \textbf{פתוחה}.
אנחנו יודעים שה-\L{$A_{n}$}-ים הם קבוצות פתוחות, וראינו שכל \L{$C_{n}$}
כזה הוא חיתוך של מספר \textbf{סופי} של קבוצות פתוחות. האם חיתוך סופי
של קבוצות פתוחות הוא בהכרח קבוצה פתוחה? במרחבים מטריים זה נכון ואוכיח
את זה עכשיו. ובאופן כללי? באופן כללי זו הולכת להיות \textbf{אחת מהאקסיומות}
שאנחנו נדרוש מקבוצות פתוחות. הנה, הגענו כבר אל אחת מאבני הבסיס היסודיות
שלנו!

מספיק להוכיח שחיתוך של שתי קבוצות פתוחות הוא פתוח; מכאן אפשר באינדוקציה
להרחיב את זה לחיתוך של מספר סופי כלשהו של קבוצות )בדיוק עם הגדרה של
\L{$C$}-ים כמו שעשיתי קודם(. אם כן, נניח ש-\L{$A_{1},A_{2}$} הן
קבוצות פתוחות ונוכיח ש-\L{$A_{1}\cap A_{2}$} היא פתוחה. ניקח \L{$a\in A_{1}\cap A_{2}$}
ונמצא \L{$r>0$} כך ש-\L{$B\left(a,r\right)\subseteq A_{1}\cap A_{2}$}.
עכשיו, בגלל ש-\L{$A_{1},A_{2}$} פתוחות ו-\L{$a\in A_{1},A_{2}$}
קיימים \L{$r_{1},r_{2}>0$} כך ש-\L{$B\left(a,r_{i}\right)\subseteq A_{i}$}.
נגדיר \L{$r=\min\left\{ r_{1},r_{2}\right\} $} ונקבל ש-\L{$B\left(a,r\right)\subseteq A_{i}$}
עבור \L{$i=1,2$}, כמו שרצינו. מה שהסתמכנו עליו כאן שיש במרחבים מטריים
ואין באופן כללי הוא שאם יש לנו שני כדורים אפשר תמיד לבחור את \textquotedblleft הקטן
יותר\textquotedblright{} והוא יהיה מוכל בשני. באופן כללי זה לא יהיה
בהכרח נכון.

אוקיי, אז דיברנו הרבה על נקודות הצטברות, אבל בעצם לא ראינו עדיין שום
דוגמא. הנה דוגמא פשוטה: נסתכל על הקטע הפתוח \L{$\left(0,1\right)$}.
מי נקודות ההצטברות בו? ראשית, כל הנקודות ששייכות לקטע, אבל מי עוד?
אם יש לנו סדרה של נקודות מתוך הקטע, כל הנקודות הללו חסומות מלעיל על
ידי {\beginL 1\endL} ומלרע על ידי {\beginL 0\endL} ולכן גם הגבול שלהן
חסום על ידי המספרים הללו. אז המספרים היחידים שיכולים להיות גבולות
של נקודות מהקטע והם לא בעצמם בקטע הם \L{$0,1$}, שתי נקודות הקצה.
ושתי הנקודות הללו הן באמת נקודות הצטברות.

זו דרך מועילה לחשוב על זה באופן כללי: נקודות הצטברות של קבוצה יכללו
את הנקודות שבתוך הקבוצה, ונקודות שהן ב\textquotedblright שפה\textquotedblright{}
של הקבוצה. מה זו שפה? אפשר להגדיר את זה פורמלית די בקלות. בהינתן \L{$A\subseteq X$},
נגדיר את השפה שלה \L{$\partial A$} בתור כל הנקודות \L{$x\in X$}
כך שבכל סביבה של \L{$x$} יש נקודה מתוך \L{$A$} ונקודה מתוך \L{$X\backslash A$},
המשלים של \L{$A$}. יודעים מה, אני לא אוהב לסמן משלים בצורה כזו, בואו
מעכשיו נשתמש בסימון \L{$A^{c}$} כדי לסמן את המשלים של \L{$A$} ביחס
ל-\L{$X$}.

עכשיו, מה קורה אם לוקחים את הקטע הפתוח \L{$\left(0,1\right)$} ומוסיפים
לו את השפה שלו, כלומר את הנקודות \L{$0,1$}? מקבלים את \textbf{הקטע
הסגור} \L{$\left[0,1\right]$}. אפשר להשתמש בזה בשביל הגדרה כללית:
נגיד שקבוצה היא \textbf{סגורה} אם היא כוללת את השפה שלה, \L{$\partial A\subseteq A$}.
עכשיו, לא יודע מה איתכם, אבל לי מפריע משהו בהגדרה הזו - היא נראית,
אה... שרירותית מדי? באופן כללי במתמטיקה כשיש לנו הגדרה אנחנו הרבה
פעמים תוהה מה \textbf{שקול} להגדרה הזו - אילו תכונות אחרות היינו יכולים
לבחור בשביל \textbf{להגדיר} קבוצה סגורה והיינו מקבלים את אותו הדבר
בדיוק. עבורי, הדבר הראשון שהייתי רוצה להבין הוא האם היינו חייבים לדבר
על \textbf{השפה} ולא, נגיד, על \textbf{נקודות ההצטברות} של הקבוצה.
האם אפשר להגדיר קבוצה סגורה ככזו שכוללת את כל נקודות ההצטברות שלה?

ובכן, קל לראות שאם \L{$x$} \textbf{לא שייכת} ל-\L{$A$} אבל היא עדיין
נקודת הצטברות של \L{$A$}, אז היא בשפה של \L{$A$}. כי בואו ניקח סביבה
כלשהי של \L{$x$}; בסביבה הזו יש נקודה של \L{$A$}, זו ממש ההגדרה
של היות \L{$x$} נקודת הצטברות. מצד שני, \L{$x$} הוא \textbf{לא}
הנקודה הזו כי הוא לא שייך ל-\L{$A$}, ולכן באותה סביבה של \L{$x$}
יש גם נקודה ב-\L{$A^{c}$}; הנקודה \L{$x$} עצמה. הוכחנו שבכל סביבה
של \L{$x$} קיימת נקודה ב-\L{$A$} ונקודה ב-\L{$A^{c}$} ולכן \L{$x\in\partial A$}.

כלומר, נקודות ההצטברות היחידות של \L{$A$} הן א( נקודות ב-\L{$A$}
עצמה )לא בהכרח כולן!( ו-ב( נקודות ב-\L{$\partial A$}. לכן אם \L{$A$}
סגורה, היא בוודאי כוללת את כל נקודות ההצטברות שלה.

זה כיוון אחד, אבל בשביל שקילות צריך להוכיח שאם \L{$A$} קבוצה שכוללת
את כל נקודות ההצטברות שלה אז היא סגורה, כלומר היא כוללת את \L{$\partial A$}.
זה עוד יותר קל - כל נקודת שפה היא נקודת הצטברות, כי הקריטריון של נקודת
שפה \textbf{כולל} את הקריטריון של נקודת הצטברות - בכל סביבה של \L{$x$}
חייב להיות איבר מ-\L{$A$} )ובנוסף עוד משהו אבל זה לא חשוב לצורך העניין
הזה(. זה מסיים את הוכחת השקילות הזו.

אבל יש עוד קריטריון שקול שמאפשר הגדרה של קבוצות סגורות, והוא לטעמי
עוד יותר מעניין. בשביל לראות אותו, בואו נסתכל שוב על הקטע \L{$\left[0,1\right]$}.
הוא חי בתוך המרחב המטרי \L{$X=\mathbb{R}$} ואפשר לדבר על \textbf{המשלים}
שלו ביחס ל-\L{$X$}. המשלים הזה הוא הקבוצה \L{$\left(-\infty,0\right)\cup\left(1,\infty\right)$}
- איחוד של שני קטעים פתוחים )עם קצוות אינסופיים, אבל זה פחות רלוונטי
כאן(. עכשיו, איחוד של שני קטעים פתוחים הוא \textbf{קבוצה פתוחה}. ואם
נסתכל על עוד דוגמאות נראה את התבנית הזו חוזרת על עצמה - אם אנחנו לוקחים
את המשלים של קבוצה סגורה אנחנו מקבלים קבוצה פתוחה. ואם אנחנו לוקחים
את המשלים של קבוצה פתוחה, אנחנו מקבלים קבוצה סגורה. זו תוצאה שבהחלט
שווה להוכיח במפורש כדי להבין איך זה קורה.

נתחיל מלהניח ש-\L{$A$} פתוחה. עכשיו, אצלנו \textquotedblleft קבוצה
פתוחה\textquotedblright{} היא מושג שנבנה על עולם המושגים של מרחב מטרי
- קבוצה היא פתוחה אם לכל נקודה בה יש כדור פתוח סביב הנקודה הזו שמוכל
כולו בקבוצה. אז אם \L{$A$} פתוחה, איך אני אוכיח ש-\L{$A^{c}$} סגורה?
מספיק לי לקחת נקודת הצטברות של \L{$A^{c}$} ולהראות שהיא \textbf{לא}
שייכת ל-\L{$A$}. תהא \L{$x$} נקודת הצטברות שכזו של \L{$A^{c}$}
ונניח בשלילה ש-\L{$x\in A$}. זה אומר שקיים כדור פתוח \L{$B\left(x,r\right)\subseteq A$}
עם \L{$r>0$}. מה קיבלנו? הכדור הפתוח הזה הוא \textbf{סביבה} של \L{$x$},
ומצד אחד הוא מוכל כולו ב-\L{$A$} ומצד שני מכיוון ש-\L{$x$} נקודת
הצטברות, כל סביבה שלה כוללת איבר של \L{$A^{c}$} ולכן \textbf{לא מוכלת}
ב-\L{$A$}. קיבלנו סתירה להנחת השלילה שלנו ש-\L{$x\in A$} ולכן \L{$x\in A^{c}$}
וזה מה שרצינו להוכיח.

עכשיו נניח ש-\L{$A$} סגורה ונוכיח ש-\L{$A^{c}$} פתוחה. כלומר, אני
אקח \L{$a\in A^{c}$} ואני אוכיח שיש לה סביבה שכולה מוכלת ב-\L{$A^{c}$}.
במילים אחרות, שיש לה סביבה שאין בה איבר של \L{$A$}. נניח ש\textbf{בכל}
סביבה של \L{$a$} יש איבר של \L{$A$}; בנוסף לכך יש בכל סביבה כזו
את \L{$a$} עצמה שהיא איבר של \L{$A^{c}$}, ולכן על פי ההגדרה, \L{$a$}
היא \textbf{נקודת שפה} של \L{$A$}. אבל \L{$A$} סגורה, ולכן מכילה
את השפה שלה, ולכן \L{$a\in A$} - סתירה. זה מסיים גם את הכיוון הזה
של ההוכחה, שהיה פשוט להפתיע - אפילו לא היינו צריכים לדבר על כדורים
- כי בניגוד לקבוצה פתוחה שאותה הגדרנו עם כדורים שהם מושג \textbf{מטרי},
את הקבוצה הסגורה הגדרנו עם נקודות הצטברות ונקודות שפה, שהם מושגים
\textbf{טופולוגיים}.

אבל מה בעצם כתבתי למעלה? אין לי הגדרה פורמלית של מה זה מושג \textquotedblleft מטרי\textquotedblright{}
ומה זה מושג \textquotedblleft טופולוגי\textquotedblright , אבל בשלב
הזה כבר יש אינטואיציות. המושגים המטריים הם אלו שמערבים את המטריקה
באופן ישיר, ובפרט \textquotedblleft כדורים פתוחים\textquotedblright .
המושגים הטופולוגיים הם אלו שלא נזקקים למושגים המטריים ישירות - קבוצות
פתוחות, סגורות, נקודות הצטברות, נקודות שפה, גבולות... את רוב המושגים
שלנו תיארנו בלי להזדקק למטריקות. מתי כן נזקקנו ממש למטריקות? כדי להוכיח
\textbf{תכונות} של המרחב שלנו שניתן לנסח באמצעות המושגים הללו. ראינו
שתי דוגמאות: הוכחנו שמרחב מטרי הוא \textbf{האוסדורפי}, כלומר ניתן
להפריד כל זוג נקודות באמצעות קבוצות פתוחות; ההגדרה של האוסדורפיות
היא גם כן טופולוגיות, אבל ההוכחה שמרחב מטרי הוא האוסדורפי מסתמכת על
התכונות המטריות. בדומה גם ראינו תכונה אחרת, על קיום של סדרה \textbf{בת
מניה }של סביבות שהיא מה שהבטיח שכל נקודת הצטברות תהיה גבול - התכונה
הזו נקראת \textbf{אקסיומת המניה הראשונה} וגם אותה הגדרנו באופן טופולוגי
והוכחנו באופן מטרי.

בקיצור, נראה שעיקר מה שאנחנו עושים הוא טופולוגיה, והחלק המטרי הוא
כמעט זניח. אז האם אפשר... לוותר עליו? בינתיים, בלי מטריקה אין לנו
את המושג של \textbf{קבוצה פתוחה} שממנו קיבלנו הכל; אבל אולי אפשר \textbf{להתחיל}
עם קבוצות פתוחות? העולם שלנו הוא קבוצה \L{$X$}. לקבוצה \L{$X$} הזו
יש תת-קבוצות שונות ומשונות; אוסף כל תתי הקבוצות של \L{$X$} מסומן
\L{$\mathcal{P}\left(X\right)$}. אולי פשוט נתחיל לדבר על טופולוגיה
על ידי כך שנבחר אוסף של תת-קבוצות, \L{$\mathcal{T}\subseteq\mathcal{P}\left(X\right)$},
נקרא לאיברים שלו \textbf{קבוצות פתוחות} ונתקדם משם?

אבל זו כמובן הגדרה כללית \textbf{מדי}. כשמתחילים להוכיח משפטים ולהגדיר
מושגים, צריך להניח משהו על מה ש-\L{$\mathcal{T}$} מקיימת. אילו תכונות
שווה לבחור? ובכן, הזכרתי כבר אחת במסגרת ההוכחה שלמעלה - השתלם לי ש\textbf{חיתוך}
של שתי קבוצות פתוחות יהיה גם הוא קבוצה פתוחה. ואם זה עובד עבור שתי
קבוצות, אפשר באינדוקציה להרחיב את זה לחיתוך של מספר \textbf{סופי}
כלשהו של קבוצות פתוחות - אבל לא אינסוף. תחשבו למשל על המרחב המטרי
\L{$\mathbb{R}$} ועל הקבוצות הפתוחות \L{$\left(-\varepsilon,\varepsilon\right)$}
עבור \L{$\varepsilon>0$} - החיתוך של כל אינסוף הקבוצות הללו הוא \L{$\left\{ 0\right\} $},
וזו לא קבוצה פתוחה כי כל כדור פתוח סביב \L{$0$} יכיל נקודות חוץ מ-{\beginL 0\endL}
עצמה. אם אנחנו הולכים לבחור תכונות כלשהן עבור טופולוגיה, אז צריך \textbf{לכל
הפחות} שהן יתקיימו עבור המרחבים המטריים, ולכן אני הולך לדרוש רק שחיתוך
של מספר סופי של קבוצות פתוחות יהיה פתוח.

מצד שני, הקבוצה \L{$\left\{ 0\right\} $} היא בהחלט \textbf{סגורה}.
אין לה נקודות שפה חוץ מ-{\beginL 0\endL} עצמה, כי כל נקודה אחרת \L{$a$}
תהיה במרחק כלשהו מ-{\beginL 0\endL} ואז אפשר לקחת כדור פתוח סביב \L{$a$}
שהוא קטן ולא מצליח להגיע עד {\beginL 0\endL}. וחוץ מזה, אפשר להסתכל
על \textbf{המשלים} של \L{$\left\{ 0\right\} $} שהוא הקבוצה \L{$\left(-\infty,0\right)\cup\left(0,\infty\right)$}
שאנחנו יודעים שהיא פתוחה וראינו שהמשלים של קבוצה פתוח הוא קבוצה סגורה.
עכשיו, את \L{$\left\{ 0\right\} $} אפשר להציג גם בתור חיתוך אינסופי
של קבוצות \textbf{סגורות}: \L{$\left\{ 0\right\} =\bigcap_{\varepsilon>0}\left[-\varepsilon,\varepsilon\right]$}.
שימו לב שזה אינסוף \textbf{גדול}, לא בן מניה אפילו. אז האם זו תכונה
שמתקיימת באופן כללי? אם \L{$\left\{ C_{i}\right\} _{i\in\Lambda}$}
היא משפחה של קבוצות סגורות שמאונדקסות על ידי אברי הקבוצה \L{$\Lambda$}
)שאינה בהכרח בת מניה(, האם \L{$C=\bigcap C_{i}$} תהיה קבוצה סגורה?
התשובה לזה פשוטה למדי: נניח ש-\L{$x$} היא נקודת הצטברות של \L{$C$}.
ניקח קבוצה שרירותית כלשהי מהחיתוך, \L{$C_{i}$}, ונראה ש-\L{$x\in C_{i}$}.
מכיוון ש-\L{$C_{i}$} סגורה, מספיק לראות ש-\L{$x$} היא נקודת הצטברות
של \L{$C_{i}$}. אז ניקח סביבה \textbf{כלשהי} של \L{$x$}. מכיוון
ש-\L{$x$} היא נקודת הצטברות של \L{$C$}, הסביבה הזו כוללת נקודה שונה
מ-\L{$x$} מתוך \L{$C\subseteq C_{i}$} כך שהנקודה הזו שייכת גם ל-\L{$C_{i}$},
וסיימנו.

על פניו זה נראה שהוכחנו את הטענה הזו בלי להיעזר בכלום! כל מה שהשתמשתי
בו הוא ההגדרה של נקודת הצטברות, שהשתמשה בתורה במושג של \textquotedblleft סביבה\textquotedblright{}
שממנו לא דרשנו כלום חוץ מכך שיהיה קבוצה. אז אם אני מגדיר קבוצה סגורה
בתור \textquotedblleft קבוצה שכוללת את כל נקודות ההצטברות שלה\textquotedblright ,
אני מקבל מייד שחיתוך כלשהו של קבוצות סגורות הוא קבוצה סגורה. והאם
אני גם יכול אוטומטית להוכיח על קבוצות כלשהן שהן סגורות? ובכן, כן,הקבוצה
\L{$X$} של המרחב כולו היא בוודאי סגורה כי היא כוללת את \textbf{כל}
הנקודות, ולכן גם את כל נקודות ההצטברות שלה. חוץ ממנה? אני לא יכול
להגיד כלום בודאות. אפילו עבור הקבוצה הריקה \L{$\emptyset$} לא. לכאורה
אין ל-\L{$\emptyset$} נקודות הצטברות בכלל, כי קחו \L{$x$} כלשהי
וסביבה \L{$U$} כלשהי של \L{$x$}; הסביבה הזו לא כוללת אף איבר של
\L{$\emptyset$} ולכן הקריטריון לנקודת הצטברות לא מתקיים. אבל בשביל
שזה יעבוד, צריך שתהיה קיימת סביבה כלשהי של \L{$x$}. אם \L{$x$} היא
נקודה שלא שייכת לאף סביבה, אז התנאי שנדרש ממנה כדי להוות נקודת הצטברות
של \L{$\emptyset$} מתקיים \textbf{באופן ריק}.

בואו נעצור לרגע כדי לעשות סדר בים הטענות שעלו פה. הנה דרך אחת להציג
את מה שעשינו:
\begin{itemize}
\item בהינתן \L{$X$} כלשהי, לקחנו \L{$\mathcal{T}\subseteq\mathcal{P}\left(X\right)$}
וקראנו לאיברים שלה \textbf{קבוצות פתוחות}.
\item לכל \L{$x\in X$}, קראנו בשם \textbf{סביבה} של \L{$x$} לקבוצה פתוחה
\L{$U\in\mathcal{T}$} כלשהי כך ש-\L{$x\in U$}.
\item הגדרנו \textbf{נקודת הצטברות} של קבוצה \L{$A$} בתור נקודה \L{$x$}
שכל סביבה שלה כוללת איבר של \L{$A$} ששונה מ-\L{$x$}.
\item הגדרנו \textbf{קבוצה סגורה} בתור קבוצה שכוללת את כל נקודות ההצטברות
שלה.
\item ראינו שחיתוך כלשהו של קבוצות סגורות הוא קבוצה סגורה.
\item ראינו ש-\L{$X$} הוא קבוצה סגורה.
\end{itemize}
את כל זה - עשינו בלי אקסיומות והנחות כלשהן. אבל פרט לכל זה, כשהנחנו
ש-\L{$X$} הוא מרחב מטרי ו\textquotedblright קבוצה פתוחה\textquotedblright{}
היא קבוצה שמוגדרת בצורה ספציפית בעזרת המטריקה )קבוצה שלכל איבר שלה
יש \textbf{כדור פתוח} סביבו שמוכל כולו בקבוצה - \textquotedblleft כדור
פתוח\textquotedblright{} הוא מושג מטרי( ראינו עוד כל מני דברים:
\begin{itemize}
\item המשלים של קבוצה סגורה הוא קבוצה פתוחה.
\item המשלים של קבוצה פתוחה הוא קבוצה סגורה.
\item החיתוך של מספר סופי של קבוצות פתוחות הוא פתוח.
\item \L{$X$} היא קבוצה פתוחה )טוב, לא הראיתי את זה, אבל זה ברור מההגדרה(.
\end{itemize}
עכשיו בואו ננסה לשלב את זה במה שאנחנו כבר יודעים. יש תוצאה בסיסית
בתורת הקבוצות שמקשרת בין חיתוכים, איחודים ומשלימים, שנקראת \textbf{כלל
דה-מורגן}. בגרסה הבסיסית שלו, כלל דה-מורגן אומר שעבור שתי קבוצות \L{$A,B$}
כלשהן מתקיים \L{$\left(A\cap B\right)^{c}=A^{c}\cup B^{c}$} ובדומה
מתקיים \L{$\left(A\cup B\right)^{c}=A^{c}\cap B^{c}$}. כלומר, אפשר
\textquotedblleft להכניס את המשלים פנימה\textquotedblright{} במחיר
של היפוך פעולת חיתוך לפעולת איחוד, והיפוך של איחוד לחיתוך.

את הדבר הזה אפשר להוכיח באותה מידה על איחודים וחיתוכים כלליים:
\begin{itemize}
\item \L{$\left(\bigcup_{i}A_{i}\right)^{c}=\bigcap_{i}A^{c}_{i}$}
\item \L{$\left(\bigcap_{i}A_{i}\right)^{c}=\bigcup_{i}A^{c}_{i}$}
\end{itemize}
איך למשל אני מוכיח את הטענה הראשונה מבין שני אלו? אני לוקח \L{$x\in\left(\bigcup_{i}A_{i}\right)^{c}$}
. על פי ההגדרה של משלים, \L{$x\notin\bigcup_{i}A_{i}$}. אם לא שייכים
לאיחוד אסור להיות שייכים לאף אחד מאיבריו, על פי הגדרת איחוד, כך ש-\L{$\forall_{i}x\notin A_{i}$}.
במילים אחרות, \L{$\forall_{i}x\in A^{c}_{i}$} ולכן מהגדרת חיתוך,
\L{$x\in\bigcap A^{c}_{i}$}.

עכשיו אפשר להחיל את דה-מורגן על מה שראינו בסיוע הטענה שהמשלים של קבוצה
סגורה הוא קבוצה פתוחה וההפך. ניקח את הטענה \textquotedblleft חיתוך
כלשהו של קבוצות סגורות הוא קבוצה סגורה\textquotedblright . אם יש לנו
חיתוך כזה, אז המשלים שלו יהיה קבוצה פתוחה, וגם המשלים של כל אחד מהאיברים
בחיתוך הוא קבוצה פתוחה - ולכן על פי דה מורגן \textbf{איחוד} של מספר
כלשהו של קבוצות \textbf{פתוחות} הוא קבוצה פתוחה. בנוסף, אם \L{$X$}
פתוחה, אז המשלימה של \L{$\emptyset$} היא סגורה. אז מה שאנחנו רואים
שהוא שבמרחב מטרי מתקיימות שלוש התכונות הבאות:
\begin{itemize}
\item \L{$\emptyset,X\in\mathcal{T}$}
\item אם יש לנו אוסף קבוצות \L{$A_{i}\in\mathcal{T}$} אז \L{$\bigcup_{i}A_{i}\in\mathcal{T}$}
\item אם יש לנו מספר סופי של קבוצות \L{$A_{1},\ldots,A_{n}\in\mathcal{T}$}
אז \L{$\bigcap^{n}_{i=1}A_{i}\in\mathcal{T}$}
\end{itemize}
ראינו שבמרחב מטרי שלוש התכונות הללו מתקיימות; יהיה לנו נוח מכאן ואילך
לקחת את התכונות הללו בתור \textbf{אקסיומות}, כלומר בתור הנחות היסוד
שלנו על מה \L{$\mathcal{T}$} מקיימת. זו נקודת המוצא שממנה נתחיל לפתח
את הטופולוגיה, והרעיון הוא שכל מרחב שיש בו אוסף של \textquotedblleft קבוצות
פתוחות\textquotedblright{} שאפשר להוכיח שמקיימות את האקסיומות )כפי
שעשינו זאת עבור מרחבים מטריים( הולך להנות אוטומטית מכל מה שנוכל להוכיח
עליהן. אז בפוסט הבא נתחיל עם ההגדרה הזו ונראה אילו דברים בסיסיים נוכל
לבנות עליה.
\end{document}
